% =====================================================================
%  Inteligência Artificial — Trânsito em Braga
%  Proposta de Projeto  →  (mais tarde) Relatório Final
%
%  COMO USAR ESTE FICHEIRO
%  - Para passar de proposta a relatório, muda \tipodocumento abaixo
%    e descomenta as secções marcadas com  [RELATÓRIO]  no fim.
%  - Os comentários "% TODO" assinalam o que falta preencher.
% =====================================================================
\documentclass[11pt,a4paper]{article}

% ---------- Língua e codificação ----------
\usepackage[utf8]{inputenc}
\usepackage[T1]{fontenc}
\usepackage[portuguese]{babel}
\usepackage{lmodern}

% ---------- Página ----------
\usepackage[a4paper,margin=2.5cm]{geometry}
\usepackage{setspace}
\setstretch{1.15}
\setlength{\parskip}{0.5em}
\setlength{\parindent}{0pt}

% ---------- Matemática, tabelas, figuras ----------
\usepackage{amsmath,amssymb}
\usepackage{booktabs}
\usepackage{tabularx}
\usepackage{array}
\usepackage{graphicx}
\usepackage{float}
\usepackage{caption}
\usepackage{enumitem}
\usepackage{pgfgantt}
\setlist{itemsep=0.2em, topsep=0.3em}

% ---------- Código (para o relatório) ----------
\usepackage{xcolor}
\usepackage{listings}
\lstset{
  language=Python,
  basicstyle=\ttfamily\small,
  keywordstyle=\color{blue!60!black}\bfseries,
  commentstyle=\color{gray}\itshape,
  stringstyle=\color{green!40!black},
  frame=single, breaklines=true, columns=fullflexible,
  literate={ã}{{\~a}}1 {á}{{\'a}}1 {à}{{\`a}}1 {â}{{\^a}}1 {é}{{\'e}}1
           {ê}{{\^e}}1 {í}{{\'i}}1 {ó}{{\'o}}1 {õ}{{\~o}}1 {ô}{{\^o}}1
           {ú}{{\'u}}1 {ç}{{\c c}}1
}

% ---------- Ligações (sempre por último) ----------
\usepackage[hidelinks]{hyperref}

% ---------- Atalhos ----------
\newcommand{\tipodocumento}{Proposta de Projeto}   % mudar para: Relatório de Projeto
\newcommand{\AStar}{A$^{*}$}
\newcommand{\IDAStar}{IDA$^{*}$}
\newcolumntype{L}{>{\raggedright\arraybackslash}X}

% =====================================================================
\begin{document}

% ---------------------------------------------------------------------
%  CAPA
% ---------------------------------------------------------------------
\hypersetup{pageanchor=false}
\begin{titlepage}
  \centering
  \vspace*{\fill}

  {\LARGE\bfseries Inteligência Artificial: Trânsito em Braga\par}
  \vspace{1.2em}
  {\large \tipodocumento\par}
  \vspace{0.8em}
  {\large\itshape Chegar a horas em Braga: procura informada e cadeias de Markov na rede dos TUB\par}

  \vspace*{\fill}
  \vspace*{\fill}

  {\large\bfseries Miguel Vila\par}
  \vspace{0.6em}
  {\small Licenciatura em Ciência de Dados Aplicada\\
          Universidade Católica Portuguesa --- Braga\\[0.3em]
          \today\par}
  \vspace{2cm}
\end{titlepage}
\hypersetup{pageanchor=true}

% ---------------------------------------------------------------------
%  ÍNDICE
% ---------------------------------------------------------------------
\tableofcontents
\newpage

% =====================================================================
\section{Proposta do Projeto}
\label{sec:proposta}
% =====================================================================

\subsection{O problema}

Quem anda de autocarro em Braga sabe que o horário é mais uma indicação
do que uma garantia. Os TUB disponibilizam os horários de toda a rede em
formato GTFS, e são esses dados que o Google Maps ou o Moovit usam para
sugerir rotas. O problema é que estas aplicações calculam a rota mais
rápida assumindo que todos os autocarros passam à hora certa.

Isto falha sobretudo quando a viagem tem transbordos. Se o primeiro
autocarro chegar três ou quatro minutos atrasado, podemos perder o
segundo e acabar por demorar o dobro do previsto. Uma rota que no papel
é a mais rápida pode ser, na prática, das menos fiáveis.

A pergunta de partida do projeto é então:

\begin{quote}
\itshape A que horas devo sair, e por que rota, para chegar à
Universidade Católica antes das 9h00 com 95\,\% de certeza?
\end{quote}

\subsection{Utilidade}

O caso mais óbvio é o de qualquer estudante ou trabalhador que tem uma
hora fixa para chegar. Toda a gente acaba por sair ``um bocado mais
cedo, por precaução'', mas sem saber bem quanto. Com este trabalho
queremos pôr um número nessa margem.

Os resultados também podem interessar ao próprio operador, porque
permitem identificar os transbordos mais frágeis, em que um pequeno
atraso tem um efeito grande no tempo total de viagem.

No contexto da unidade curricular, o projeto serve para ver na prática
uma coisa que até agora só vimos em teoria: o que acontece a um problema
de procura quando o ambiente deixa de ser determinista.

\subsection{Objetivos}

Dividimos o trabalho em três objetivos, cada um a partir do anterior:

\begin{enumerate}
  \item Encontrar a rota mais rápida segundo o horário publicado,
        formulando a rede como um problema de procura e resolvendo-o com
        custo uniforme, \AStar{} e algumas variantes. Queremos também
        provar que a heurística usada é admissível.
  \item Incluir os atrasos através de uma cadeia de Markov e procurar a
        rota com menor tempo \emph{esperado}.
  \item Encontrar a rota com maior probabilidade de chegar antes de uma
        dada hora, avaliando as rotas candidatas por simulação.
\end{enumerate}

Se houver tempo, gostávamos ainda de experimentar uma versão em que o
plano se ajusta durante a viagem (por exemplo, mudar de linha no
transbordo se o primeiro autocarro vier atrasado). Isso já obriga a
formular o problema como um MDP.

\subsection{Estratégia}

O primeiro passo é tratar os dados GTFS. Para o problema não ficar
demasiado grande, vamos limitar-nos a um dia útil, a uma janela horária
de manhã (7h00--10h00) e a uma zona da cidade que inclua a UCP.

A partir daí construímos um grafo em que cada estado é um par
(paragem, hora). As ligações entre estados correspondem a viajar num
autocarro, caminhar até uma paragem próxima ou esperar. Sobre este
grafo aplicamos os algoritmos de procura dados nas aulas e comparamos
o número de nós expandidos, o tempo de execução e o custo das soluções.

Para os atrasos, definimos uma cadeia de Markov com poucos estados
(a tempo, atraso pequeno, atraso grande) e estimamos as probabilidades
de transição. Por fim, geramos várias rotas candidatas com \AStar{} e
simulamos milhares de viagens em cada uma, para ver com que frequência
se chega a horas. O resultado que procuramos são casos concretos em que
a rota mais rápida no horário não é a mais segura.

\subsubsection*{Calendário}

O plano de trabalho até à apresentação está na
Figura~\ref{fig:gantt}. Marcámos também o Teste 2 e a entrega do TPC 2,
porque calham a meio do projeto e condicionam o ritmo dessas semanas.

\begin{figure}[H]
\centering
\begin{ganttchart}[
    hgrid,
    vgrid={*6{draw=none},dotted},
    time slot format=isodate,
    x unit=0.36cm,
    y unit chart=0.55cm,
    y unit title=0.5cm,
    title height=1,
    title label font=\scriptsize,
    bar label font=\small,
    milestone label font=\small\itshape,
    bar height=0.5,
    bar/.append style={fill=blue!35, draw=blue!60!black},
    group/.append style={fill=black!70},
    milestone/.append style={fill=red!70!black, draw=red!70!black},
    milestone height=0.5,
    milestone left shift=0.1, milestone right shift=-0.1
  ]{2026-11-14}{2026-12-16}
  \gantttitle{Novembro}{17} \gantttitle{Dezembro}{16} \\
  \gantttitlelist{14,...,30}{1} \gantttitlelist{1,...,16}{1} \\
  \ganttbar{Proposta}{2026-11-14}{2026-11-18} \\
  \ganttmilestone{Entrega da proposta}{2026-11-18} \\
  \ganttbar{Dados GTFS e grafo}{2026-11-19}{2026-11-25} \\
  \ganttbar{Recolha de atrasos}{2026-11-19}{2026-12-04} \\
  \ganttbar{Procura determinística}{2026-11-26}{2026-12-01} \\
  \ganttbar{Modelo de Markov}{2026-12-02}{2026-12-06} \\
  \ganttmilestone{Teste 2}{2026-12-07} \\
  \ganttbar{Procura + Markov}{2026-12-08}{2026-12-12} \\
  \ganttmilestone{Entrega TPC 2}{2026-12-09} \\
  \ganttbar{Relatório e apresentação}{2026-12-13}{2026-12-16} \\
  \ganttmilestone{Apresentação}{2026-12-16}
\end{ganttchart}
\caption{Calendário previsto do projeto.}
\label{fig:gantt}
\end{figure}

\subsubsection*{Riscos e limitações}

A maior dificuldade é que o GTFS só tem os horários previstos, não os
atrasos reais. A ideia é registar nós próprios os atrasos em algumas
viagens e, em paralelo, perguntar aos TUB se têm dados que possam
partilhar. Se nada disso chegar, usamos probabilidades assumidas e
analisamos o quanto os resultados mudam quando essas probabilidades
mudam. O outro risco é o grafo crescer demais, o que se resolve
reduzindo a zona ou a janela horária.

% =====================================================================
\section{Descrição PEAS e propriedades do ambiente}
\label{sec:peas}
% =====================================================================

\subsection{Descrição PEAS}

O agente que estamos a desenhar é um planeador de viagens: recebe uma
origem, um destino e uma hora limite e devolve o percurso a fazer.

\begin{table}[H]
\centering
\begin{tabularx}{\textwidth}{@{}l L@{}}
\toprule
\textbf{Componente} & \textbf{Descrição} \\
\midrule
Desempenho (P) &
  Probabilidade de chegar antes da hora limite, tempo de viagem,
  número de transbordos e distância feita a pé. \\[0.3em]
Ambiente (E) &
  Rede dos TUB (paragens, linhas e horários), ruas entre paragens,
  atrasos dos autocarros e hora do dia. \\[0.3em]
Atuadores (A) &
  Indicar a próxima ação: apanhar a linha $X$ na paragem $Y$,
  ir a pé até outra paragem, esperar ou sair. \\[0.3em]
Sensores (S) &
  Ficheiros GTFS, posição e hora atuais e, na versão adaptativa,
  o atraso observado do autocarro. \\
\bottomrule
\end{tabularx}
\caption{Descrição PEAS do planeador de viagens.}
\end{table}

\subsection{Propriedades do ambiente}

Classificámos o ambiente de duas maneiras: considerando só o horário
publicado e considerando os atrasos. A diferença entre as duas colunas
é, no fundo, a razão de ser do projeto.

\begin{table}[H]
\centering
\begin{tabularx}{\textwidth}{@{}l L L@{}}
\toprule
\textbf{Propriedade} & \textbf{Só horário} & \textbf{Com atrasos} \\
\midrule
Observabilidade        & Totalmente observável & Parcialmente observável (não se sabe o atraso futuro) \\
Agentes                & Único & Único (o resto do trânsito faz parte do ambiente) \\
Determinismo           & Determinista & Estocástico \\
Episódico/sequencial   & Sequencial & Sequencial \\
Estático/dinâmico      & Estático & Dinâmico \\
Discreto/contínuo      & Discreto & Discreto \\
Conhecido/desconhecido & Conhecido & Parcialmente conhecido (probabilidades estimadas) \\
\bottomrule
\end{tabularx}
\caption{Classificação do ambiente nas sete propriedades.}
\end{table}

Com o horário apenas, o ambiente é totalmente observável, determinista,
estático e conhecido, e por isso basta procurar uma sequência de ações
antes de sair de casa. Com atrasos isso deixa de ser suficiente, e
passa a fazer sentido trabalhar com probabilidades e, numa fase mais
avançada, com políticas em vez de planos fixos.

% =====================================================================
\section{Ferramentas}
\label{sec:ferramentas}
% =====================================================================

\subsection{Dados}

Os dados principais vêm do GTFS dos TUB \cite{tubgtfs}. Os ficheiros que
vamos usar são o \texttt{stops.txt} (paragens e coordenadas), o
\texttt{routes.txt} (linhas), o \texttt{trips.txt} (viagens), o
\texttt{stop\_times.txt} (hora de passagem em cada paragem) e o
\texttt{calendar.txt} (dias em que cada serviço funciona). Os atrasos
serão registados à parte, anotando a hora prevista e a hora real de
passagem em viagens feitas por nós.
% TODO: descrever como foi feita a recolha

\subsection{Formulação como problema de procura}

Seguindo os seis componentes vistos nas aulas:

\begin{itemize}
  \item estados: pares $(p, t)$, em que $p$ é uma paragem e $t$ a hora;
  \item estado inicial: $(p_{\text{origem}}, t_{\text{saída}})$;
  \item objetivo: estar na paragem de destino, a qualquer hora $t$;
  \item ações: apanhar um autocarro que passe em $p$ depois de $t$,
        caminhar até uma paragem próxima ou esperar;
  \item modelo de transição: dado pelos horários do GTFS e, a pé, por uma
        velocidade média de caminhada;
  \item custo: o tempo que passa, com uma possível penalização por
        cada transbordo.
\end{itemize}

\subsection{Algoritmos de procura}

Vamos comparar os algoritmos da Tabela~\ref{tab:algoritmos} usando as
quatro medidas do capítulo 3 de \cite{aima}: completude, otimalidade,
tempo (nós expandidos) e memória (tamanho máximo da fronteira).

\begin{table}[H]
\centering
\begin{tabularx}{\textwidth}{@{}l l L@{}}
\toprule
\textbf{Algoritmo} & \textbf{Avaliação} & \textbf{Para que serve aqui} \\
\midrule
Custo uniforme      & $f = g$                 & Solução ótima de referência \\
\AStar{}            & $f = g + h$             & Algoritmo principal; ótimo se $h$ for admissível \\
\AStar{} pesado     & $f = g + W\,h$, $W>1$   & Ver quanto se ganha em rapidez e quanto se perde no custo ($\le W \cdot C^{*}$) \\
Procura em feixe    & $k$ melhores nós        & Memória limitada, sem garantias \\
\IDAStar{}          & limite sobre $f$        & Memória linear (opcional) \\
\bottomrule
\end{tabularx}
\caption{Algoritmos de procura a comparar.}
\label{tab:algoritmos}
\end{table}

\subsection{Heurística}

Como heurística usamos o tempo que se levaria a chegar ao destino em
linha reta, à maior velocidade possível na rede:
\[
  h(p, t) = \frac{d(p, p_{\text{destino}})}{v_{\max}},
\]
onde $d$ é a distância em linha reta e $v_{\max}$ é a maior velocidade
média registada num troço de autocarro, calculada a partir do próprio
GTFS.

Esta heurística é admissível porque não há nenhuma forma de percorrer a
distância $d$ em menos de $d / v_{\max}$, uma vez que nenhum autocarro
(e muito menos alguém a pé) vai mais depressa do que $v_{\max}$. Logo,
$h(n) \le h^{*}(n)$ para qualquer nó $n$.

É também consistente. Uma ação que leva de $p$ a $p'$ demora pelo menos
$d(p,p')/v_{\max}$ e, pela desigualdade triangular,
$d(p,p_{\text{destino}}) \le d(p,p') + d(p',p_{\text{destino}})$. Daqui
resulta $h(n) \le c(n,n') + h(n')$.

% TODO (relatório): comparar com h = 0 e com outras heurísticas;
% calcular o fator de ramificação efetivo b*.

\subsection{Cadeias de Markov}

Para modelar os atrasos usamos uma cadeia de Markov com três estados:
\[
  S = \{\textsf{a tempo},\ \textsf{atraso pequeno},\ \textsf{atraso grande}\}.
\]
A ideia é que o atraso de um autocarro no troço seguinte só depende do
atraso que traz no troço atual:
\[
  P(X_{k+1} = j \mid X_k = i, X_{k-1}, \dots, X_0) = P(X_{k+1} = j \mid X_k = i) = p_{ij}.
\]

As probabilidades estimam-se contando transições nos dados recolhidos,
$\hat{p}_{ij} = n_{ij} / \sum_{j} n_{ij}$. Se houver dados suficientes,
fazemos uma matriz para a hora de ponta e outra para o resto do dia. A
partir da matriz $P$ podemos ainda calcular a distribuição estacionária
($\pi = \pi P$), que indica a proporção de tempo que um autocarro passa
em cada estado, e a distribuição do atraso ao fim de $k$ troços,
$\mu_0 P^{k}$.

\subsection{Simulação de Monte Carlo}

A probabilidade de chegar a horas não se vai somando ao longo do
caminho, ao contrário do tempo, por isso não pode ser usada diretamente
como custo no \AStar{}. A solução é avaliar cada rota candidata por
simulação: em cada viagem simulada sorteamos os atrasos a partir da
cadeia de Markov, verificamos se os transbordos se conseguem fazer e
registamos a hora de chegada. A fração de simulações que chegam a tempo
dá-nos uma estimativa da probabilidade.

\subsection{Processos de decisão de Markov (extensão)}

Um MDP é definido por $(S, A, P, R)$: estados, ações, probabilidades de
transição $P(s' \mid s, a)$ e recompensas. A diferença em relação à
procura é que a solução deixa de ser uma sequência de ações e passa a
ser uma política $\pi(s)$, que indica o que fazer em cada estado. No
nosso caso, o estado teria a paragem, a hora e o atraso observado, o que
permitiria decisões do tipo ``se o autocarro vier atrasado, troco de
linha no transbordo''.

\subsection{Software}

O projeto será feito em Python, com \texttt{pandas} e \texttt{numpy}
para os dados e as matrizes, a biblioteca \texttt{aima-python}
\cite{aimapython} para a estrutura do problema e os algoritmos de
procura, \texttt{matplotlib} para os gráficos e \texttt{folium} para
desenhar a rede e as rotas num mapa.

% =====================================================================
%  [RELATÓRIO] Secções a descomentar quando o projeto avançar
% =====================================================================

% \section{Dados e construção do grafo}
% \label{sec:dados}
% % TODO: zona e janela horária escolhidas, número de paragens/arestas,
% %       figura do mapa da rede (folium -> captura de ecrã).
% % \begin{figure}[H]
% %   \centering
% %   \includegraphics[width=0.85\textwidth]{figuras/mapa_rede.png}
% %   \caption{Rede filtrada dos TUB.}
% % \end{figure}

% \section{Procura determinística}
% \label{sec:procura}
% % TODO: implementação, pares origem/destino testados,
% %       tabela de nós expandidos / tempo / custo para cada algoritmo,
% %       gráfico do "eixo W" (W = 0, 1, 1.5, 2).

% \section{Modelo de Markov dos atrasos}
% \label{sec:markov}
% % TODO: recolha de dados, matriz de transição estimada,
% %       distribuição estacionária, validação por simulação.

% \section{Resultados: horário vs.\ tempo esperado vs.\ probabilidade}
% \label{sec:resultados}
% % TODO: comparação dos três critérios, exemplos em que a rota mais
% %       rápida não é a mais fiável, gráfico P(chegar antes das 9h00)
% %       em função da hora de saída.

% \section{Conclusões e trabalho futuro}
% \label{sec:conclusoes}
% % TODO: o que se aprendeu, limitações, extensão para MDP.

% =====================================================================
\begin{thebibliography}{9}

\bibitem{aima}
S.~Russell e P.~Norvig.
\newblock \emph{Artificial Intelligence: A Modern Approach}.
\newblock 4.ª ed., Global Edition, Pearson, 2021.

\bibitem{tubgtfs}
Transportes Urbanos de Braga.
\newblock \emph{GTFS --- Transportes Urbanos de Braga}.
\newblock \url{https://dados.gov.pt/pt/datasets/gtfs-transportes-urbanos-de-braga/}.

\bibitem{aimapython}
AIMA Code.
\newblock \emph{aima-python: Python implementation of algorithms from AIMA}.
\newblock \url{https://github.com/aimacode/aima-python}.

\bibitem{gtfsref}
MobilityData.
\newblock \emph{GTFS Schedule Reference}.
\newblock \url{https://gtfs.org/schedule/reference/}.

\end{thebibliography}

\end{document}
